%============================== INTRODUCTION =================================
\chapter{Introduction générale}

\section{Contexte : la publicité digitale}

La publicité digitale s'est imposée, en quelques années, comme un levier
central des stratégies de communication des annonceurs. Elle repose sur des
plateformes technologiques capables de diffuser des messages ciblés auprès
d'audiences précises, puis de mesurer finement leurs performances. Dans cet
écosystème, les \emph{Demand-Side Platforms} (DSP) occupent une place
particulière : ce sont les plateformes côté \og demande \fg{} qui permettent
aux annonceurs (ou aux agences qui les représentent) de planifier, configurer,
diffuser et suivre leurs campagnes publicitaires depuis une interface unique.

\section{Les campagnes \emph{drive-to-store}}

Parmi les objectifs publicitaires, le \textbf{drive-to-store} vise un résultat
très concret : transformer une exposition publicitaire en ligne en une
\textbf{visite physique dans un point de vente}. Ce type de campagne est
particulièrement pertinent pour les enseignes disposant d'un réseau de
magasins (grande distribution, banques, hard-discount, etc.), qui cherchent à
générer du trafic en boutique plutôt que de simples clics.

Une campagne \emph{drive-to-store} mobilise plusieurs ingrédients : un
\textbf{ciblage géographique} autour des magasins (le \emph{geofencing}), des
\textbf{créatives adaptées} à chaque point de vente (nom, adresse, horaires),
et un \textbf{reporting} orienté sur la fréquentation et la performance locale.
La coordination de ces éléments justifie l'existence d'une plateforme dédiée.

\section{Importance d'une plateforme de gestion de campagnes}

Gérer manuellement l'ensemble de ces étapes --- création de campagne, import
des magasins, définition des zones de diffusion, production des variantes
créatives, suivi des indicateurs --- serait fastidieux et source d'erreurs.
Une \textbf{plateforme centralisée} apporte plusieurs bénéfices :

\begin{itemize}
  \item un \textbf{parcours guidé} et cohérent, de la création jusqu'au
        reporting ;
  \item la \textbf{réutilisation} des données (magasins, annonceurs) d'un
        module à l'autre ;
  \item une \textbf{maîtrise des accès} par rôle (administration, achat média,
        lecture seule) ;
  \item une \textbf{visualisation} claire des performances et des zones de
        diffusion.
\end{itemize}

\section{Motivation et périmètre du projet}

Le présent projet consiste à concevoir et développer, pour SBS Data Factory,
un \textbf{prototype de plateforme Drive-to-Store DSP} illustrant ce parcours
de bout en bout. Il s'agit d'une application web fullstack, et non d'un
\emph{bidder} temps réel (RTB) ni d'un backend OpenRTB : l'accent est mis sur
l'\textbf{expérience de gestion de campagne} et sur une \textbf{architecture
logicielle propre et évolutive}, capable d'accueillir ultérieurement une base
de données réelle et une authentification sécurisée.

Ce rapport présente successivement l'organisme d'accueil, la problématique et
les objectifs, la méthodologie de travail, l'analyse fonctionnelle,
l'architecture technique, la conception de la base de données, la réalisation
module par module, les tests et validations, la documentation et la livraison,
puis les limites et perspectives, avant de conclure. Un souci constant de
\textbf{transparence} guide l'ensemble : les parties simulées ou \emph{mockées}
du prototype sont explicitement signalées.
